\documentclass[12pt,a4paper]{article}

% ------------------------------------------------
% Pacotes
% ------------------------------------------------
\usepackage[utf8]{inputenc}
\usepackage[T1]{fontenc}
\usepackage[brazil]{babel}
\usepackage{graphicx}
\usepackage[table]{xcolor}
\usepackage{array}
\usepackage{tabularx}
\usepackage{geometry}
\usepackage{fancyhdr}
\usepackage{mathptmx} % Times New Roman
\usepackage{enumitem}\usepackage{amsfonts}
\usepackage{amsmath} 
\usepackage{amssymb}
\usepackage{tikz}
\usetikzlibrary{positioning, shapes.geometric, calc}

% TikZ Styles for ER Diagrams
\tikzset{
    block/.style={rectangle, draw, thick, minimum width=2cm, minimum height=1cm, fill=white},
    diamond_node/.style={diamond, draw, thick, minimum width=2cm, minimum height=1cm, fill=white, aspect=2},
    attribute/.style={ellipse, draw, thin, minimum width=1.5cm, minimum height=0.6cm, fill=white, font=\scriptsize},
    line/.style={draw, thick},
}

% ------------------------------------------------ 
% Configuração de página
% ------------------------------------------------
\geometry{
    a4paper,
    left=2cm,
    right=2cm,
    top=2cm,
    bottom=2cm,
    headheight=15pt 
}

% Caixa para resposta discursiva
\newcommand{\resposta}[1]{
\noindent
\framebox[\textwidth]{
\begin{minipage}{0.98\textwidth}
\vspace{#1}
\end{minipage}
}
}

\pagestyle{fancy}
\fancyhf{}
\renewcommand{\headrulewidth}{0pt}

\setlength{\parindent}{0pt}

% ------------------------------------------------
% Comandos auxiliares
% ------------------------------------------------
\newcolumntype{L}[1]{>{\raggedright\arraybackslash}p{#1}}
\newcolumntype{C}[1]{>{\centering\arraybackslash}p{#1}}

% ------------------------------------------------
% Documento
% ------------------------------------------------
\begin{document}

% ------------------------------------------------
% Cabeçalho
% ------------------------------------------------
\begin{tabularx}{\textwidth}{@{} L{4cm} X @{}}
    \raisebox{-0.5cm}{\includegraphics[width=4cm]{Universidade_de_Rio_Verde.jpg}} &
    \hfill
    \begin{tabular}{|>{\centering\arraybackslash}m{10cm}|>{\raggedright\arraybackslash}p{1.8cm}|}
        \hline
        \cellcolor{black}\color{white}\rule{0pt}{1.2cm}\textbf{LISTA DE REVISÃO - N2} 
        & 
        \scriptsize Nota \tabularnewline
        \hline
    \end{tabular} 
\end{tabularx}
\vspace{0.3cm}

% ------------------------------------------------  
% Tabela principal
% ------------------------------------------------
\begin{tabularx}{\textwidth}{|X|L{2cm}|}
\hline
\textbf{Curso} & \textbf{Campus} \\
\textit{Engenharia de Software} & \textit{Rio Verde} \\
\hline
\multicolumn{2}{|l|}{\textbf{Disciplina}} \\  
\multicolumn{2}{|l|}{ \textit{Banco de dados}} \\
\hline
\textbf{Nome do(a) acadêmico(a)} & \textbf{Valor} \\
& --- \\
\hline
\textbf{Professor} & \textbf{Turma} \\
\textit{ Sergio Souza Novak} & \textit{B} \\
\hline
\textbf{Data de Entrega} & \textbf{Código} \\
\textit{--/05/2026} & \textit{ESW417} \\
\hline
\end{tabularx}
 
\vspace{0.5cm}

\begin{center}
    \rule{\textwidth}{0.4pt}
    \textbf{QUESTÕES TEÓRICAS (NORMALIZAÇÃO E CONCEITOS)}
\end{center}

\begin{enumerate}
     \item \textbf{Os diagramas de Venn abaixo representam visualmente diferentes operações de JOIN entre duas tabelas, A (esquerda) e B (direita). Analise as áreas sombreadas e assinale a alternativa que identifica corretamente os tipos de JOIN representados nas Figuras I, II e III, respectivamente:}
    
    \begin{center}
        \begin{tikzpicture}[thick, scale=0.8]
            % Figura I: RIGHT JOIN
            \begin{scope}
                \fill[gray!40] (1.2,0) circle (1cm);
            \end{scope}
            \draw (0,0) circle (1cm) node[xshift=-0.5cm] {A};
            \draw (1.2,0) circle (1cm) node[xshift=0.5cm] {B};
            \node at (0.6, -1.5) {\textbf{Figura I}};
        \end{tikzpicture}
        \hspace{1cm}
        \begin{tikzpicture}[thick, scale=0.8]
            % Figura II: FULL OUTER JOIN
            \fill[gray!40] (0,0) circle (1cm);
            \fill[gray!40] (1.2,0) circle (1cm);
            \draw (0,0) circle (1cm) node[xshift=-0.5cm] {A};
            \draw (1.2,0) circle (1cm) node[xshift=0.5cm] {B};
            \node at (0.6, -1.5) {\textbf{Figura II}};
        \end{tikzpicture}
        \hspace{1cm}
        \begin{tikzpicture}[thick, scale=0.8]
            % Figura III: LEFT JOIN
            \fill[gray!40] (0,0) circle (1cm);
            \begin{scope}
                \clip (0,0) circle (1cm);
                \fill[gray!40] (1.2,0) circle (1cm);
            \end{scope}
            \draw (0,0) circle (1cm) node[xshift=-0.5cm] {A};
            \draw (1.2,0) circle (1cm) node[xshift=0.5cm] {B};
            \node at (0.6, -1.5) {\textbf{Figura III}};
        \end{tikzpicture}
    \end{center}

    \begin{enumerate}[label=\Alph*)]
        \item I: LEFT JOIN; II: FULL JOIN; III: RIGHT JOIN.
        \item I: RIGHT JOIN; II: FULL JOIN; III: LEFT JOIN.
        \item I: RIGHT JOIN; II: INNER JOIN; III: LEFT JOIN.
        \item I: JOIN; II: FULL JOIN; III: LEFT JOIN.
        \item I: FULL JOIN; II: RIGHT JOIN; III: LEFT JOIN.
    \end{enumerate}

    \item \textbf{Escolha uma das figuras acima (I, II ou III) e escreva a sintaxe SQL básica que realizaria tal operação entre as tabelas \texttt{Clientes} (A) e \texttt{Pedidos} (B), utilizando o campo \texttt{id} como critério de junção. (Exemplo: para um \texttt{INNER JOIN}, a sintaxe seria \texttt{SELECT * FROM A INNER JOIN B ON A.id = B.id})}\\
    \resposta{2.5cm}
    \item \textbf{Sobre os conceitos fundamentais de Banco de Dados Relacionais, analise as afirmações abaixo:}
    \begin{enumerate}[label=\Roman*.]
        \item Uma Chave Estrangeira (FK) pode aceitar valores nulos, a menos que faça parte da chave primária da tabela onde está inserida.
        \item A dependência funcional transitiva ocorre quando um atributo não-chave depende de outro atributo não-chave, que por sua vez depende da chave primária.
        \item Para uma tabela estar na Segunda Forma Normal (2FN), ela deve obrigatoriamente estar na 1FN e não possuir dependências parciais.
        \item A Terceira Forma Normal (3FN) elimina redundâncias causadas por atributos multivalorados.
    \end{enumerate}
    Estão corretas apenas:
    \begin{enumerate}[label=\Alph*)]
        \item I e II.
        \item II e III.
        \item I, II e III.
        \item I, III e IV.
        \item Todas estão corretas.
    \end{enumerate}

    \item \textbf{Considere a tabela abaixo que armazena informações de Projetos e seus respectivos participantes. A tabela viola a 1FN.}
    
    \begin{center}
        \begin{tabularx}{\textwidth}{|c|c|X|}
            \hline
            \cellcolor{gray!20}\textbf{Cod\_Proj} & \cellcolor{gray!20}\textbf{Nome\_Projeto} & \cellcolor{gray!20}\textbf{Equipe (Membros)} \\ \hline
            P01 & Alpha & Ricardo (Dev), Julia (Design) \\ \hline
            P02 & Beta & Marcos (QA) \\ \hline
            P03 & Gamma & Paulo (Dev), Silvia (Dev), Ana (PO) \\ \hline
            P04 & Delta & Fabio (Front), Carla (UX) \\ \hline
            P05 & Epsilon & Renato (Back) \\ \hline
            P06 & Zeta & Lucas (Dev), Maria (Dev), Jose (PM) \\ \hline
            P07 & Eta & Bruno (Dev), Alice (Design) \\ \hline
            P08 & Theta & Igor (QA), Gabi (Dev) \\ \hline
            P09 & Iota & Samuel (Dev), Debora (Dev) \\ \hline
        \end{tabularx}
    \end{center}
    \begin{enumerate}[label=(\Alph*)]
        \item Explique por que a tabela viola a Primeira Forma Normal (1FN).\\
        \resposta{2cm}
        \item Apresente a tabela resultante após a normalização para a 1FN.\\
        \resposta{5cm}
    \end{enumerate}

    \item \textbf{Dada a tabela abaixo de MATRÍCULAS, identifique as dependências funcionais e determine em qual Forma Normal a tabela se encontra. Justifique sua resposta.}
    
    \textit{Chave Primária Composta: \textbf{\{Cod\_Aluno, Cod\_Curso\}}}
    
    \begin{center}
        \begin{tabularx}{\textwidth}{|c|c|c|X|c|}
            \hline
            \cellcolor{gray!20}\textbf{Cod\_Alu} & \cellcolor{gray!20}\textbf{Cod\_Cur} & \cellcolor{gray!20}\textbf{Data\_Insc} & \cellcolor{gray!20}\textbf{Nome\_Curso} & \cellcolor{gray!20}\textbf{Carga} \\ \hline
            100 & C1 & 10/02 & SQL Básico & 40h \\ \hline
            100 & C2 & 15/02 & Python & 60h \\ \hline
            200 & C1 & 11/02 & SQL Básico & 40h \\ \hline
            200 & C2 & 16/02 & Python & 60h \\ \hline
            300 & C3 & 20/02 & Java Advanced & 80h \\ \hline
            300 & C1 & 12/02 & SQL Básico & 40h \\ \hline
            400 & C4 & 25/02 & Docker & 30h \\ \hline
            400 & C2 & 17/02 & Python & 60h \\ \hline
            500 & C3 & 21/02 & Java Advanced & 80h \\ \hline
        \end{tabularx}
    \end{center}
    \resposta{4cm}

    \item \textbf{A Quarta Forma Normal (4FN) lida com Dependências Multivaloradas. Considere o cenário onde um Professor pode ministrar várias Disciplinas e orientar vários Projetos de Pesquisa. As disciplinas e os projetos são totalmente independentes.}
    
    Decomponha a tabela abaixo para atingir a 4FN:
    
    \begin{center}
        \begin{tabularx}{\textwidth}{|X|X|X|}
            \hline
            \cellcolor{gray!20}\textbf{Professor (PK)} & \cellcolor{gray!20}\textbf{Disciplina (PK)} & \cellcolor{gray!20}\textbf{Projeto\_Pesquisa (PK)} \\ \hline
            Novak & Banco de Dados & IA Generativa \\ \hline
            Novak & Banco de Dados & Redes Neurais \\ \hline
            Novak & Banco de Dados & NLP \\ \hline
            Novak & Programação & IA Generativa \\ \hline
            Novak & Programação & Redes Neurais \\ \hline
            Novak & Programação & NLP \\ \hline
            Novak & Eng. Requisitos & IA Generativa \\ \hline
            Novak & Eng. Requisitos & Redes Neurais \\ \hline
            Novak & Eng. Requisitos & NLP \\ \hline
            Novak & Algoritmos & IA Generativa \\ \hline
            Novak & Algoritmos & Redes Neurais \\ \hline
            Novak & Algoritmos & NLP \\ \hline
        \end{tabularx}
    \end{center}
    \resposta{6cm}

   
\end{enumerate}

\newpage
\begin{center}
    \rule{\textwidth}{0.4pt}
    \textbf{QUESTÕES PRÁTICAS (SQL)}
\end{center}

\textbf{Cenário:} Uma livraria deseja organizar seu estoque e vendas. O esquema lógico é:
\begin{itemize}[label={}]
    \item \textbf{tbAutor} (\underline{id\_autor}, nome\_autor, nacionalidade)
    \item \textbf{tbLivro} (\underline{id\_livro}, titulo, ano\_publicacao, preco, id\_autor)
    \item \textit{id\_autor referencia tbAutor}
    \item \textbf{tbVenda} (\underline{id\_venda}, data\_venda, id\_livro, qtd)
    \item \textit{id\_livro referencia tbLivro}
\end{itemize}

\begin{enumerate}
    \item Escreva o comando \textbf{DDL} para criar as tabelas \textbf{tbAutor} e \textbf{tbLivro}, garantindo as chaves primárias e a integridade\\
    \resposta{7cm}

    \item Insira \textbf{3 registros} na tabela \textbf{tbAutor} e \textbf{2 registros} na tabela \textbf{tbLivro}.\\
    \resposta{5cm}

    \item Escreva um comando SQL para aumentar o preço de todos os livros publicados antes de 2020 em 10\%.\\
    \resposta{3cm}

    \item Realize uma consulta (\textbf{SELECT}) que retorne o título do livro, o nome do autor e o preço, apenas para livros que custam mais de R\$ 50,00, ordenados pelo título.\\
    \resposta{4cm}

    \item Utilizando \textbf{JOIN}, liste todos os livros (título) e a quantidade vendida de cada um (da tabela \texttt{tbVenda}). Se um livro nunca foi vendido, ele ainda deve aparecer na lista com a quantidade nula ou zero.\\
    \resposta{4cm}

    \item Escreva uma consulta SQL que retorne o nome do autor, o título do livro e a data da venda para todas as vendas registradas no sistema. (Dica: conecte as três tabelas: \texttt{tbAutor}, \texttt{tbLivro} e \texttt{tbVenda}).\\
    \resposta{4cm}

    \item Utilizando \textbf{JOIN} e a função de agregação \texttt{SUM()}, calcule o valor total arrecadado por cada autor (soma de \texttt{preco * qtd}). O resultado deve exibir o nome do autor e o valor total.\\
    \resposta{4cm}
\end{enumerate}

\end{document}
